\chapter{NURBS curves and interpolation}
\label{chap:splines}

\section{One curve definition, two separate worlds}

The ordinary curve types are \api{nurbs_spline2<REAL>} and
\api{nurbs_spline3<REAL>}. Their coordinate storage is private, their definition
is checked at construction and after each setter, and their active-domain
endpoints are cached. The corresponding \api{nurbs_arr_spline2<REAL>} and
\api{nurbs_arr_spline3<REAL>} types provide a different storage contract: a
public, flattened \api{std::valarray<REAL>} that an optimizer can edit directly.
All four types describe rational B-spline curves and provide position,
derivative, tangent, curvature, length, and interpolation operations.

The suffix is part of the geometry contract. A 2D curve evaluates to
\api{point2<REAL>} and differentiates to \api{vector2<REAL>}; a 3D curve uses
the corresponding 3D types. Constructing an array-backed curve from an ordinary
curve, or converting it back, preserves that suffix. Embedding a 2D curve into
3D is an explicit \api{project(plane, spline)} operation, described in
Chapter~\ref{chap:entities}.

For a curve with $C$ control points, degree $p$, weights $w_i$, and knots $K_j$,
the complete definition satisfies
\begin{equation}
 C>p\geq1,\qquad \#w=C,\qquad \#K=C+p+1,
 \qquad w_i>0,\qquad K_j\leq K_{j+1}.
\end{equation}
Coordinates, weights, knots, and the strictly positive definition tolerance
must be finite. The knot domain used by the curve is
\begin{equation}
\boxed{\quad s\in[K_p,K_C]=[\api{s_min()},\api{s_max()}]\quad}
\label{eq:spline-domain}
\end{equation}
and must have positive length. It is the \emph{active} domain; the first and
last values in the stored knot vector need not be its bounds.

\begin{manualnote}[Native parameters]
The parameter is named $s$, but a general NURBS definition does not make $s$
equal to distance traveled. Keep the native knot units when passing $s$ to
evaluation, intersections, distances, or derivative functions. Interpolation
preserves the supplied station values, including a nonzero first station.
\end{manualnote}

\section{Construction from a complete definition}
\label{sec:spline-definition}

The following fragment establishes a running example in the native 2D world.
It uses five control points, cubic basis functions, nonunit weights, and a
native domain of $[0,5]$. Code fragments in this chapter use the umbrella header
\api{<nurbspath/nurbspath.hpp>}; standard-library facilities need their usual
headers. The namespace alias \api{np} below is used throughout the chapter.

\begin{lstlisting}[language=C++,caption={A complete ordinary 2D curve definition.}]
namespace np = nurbspath;
using real = double;
using point = np::point2<real>;

const std::vector<point> controls{
    {0, 0}, {1, 2}, {2, -1}, {3, 2}, {4, 0}};
const std::vector<real> weights{1, 1, 0.7, 1.5, 1};
const std::vector<real> knots{0, 0, 0, 0, 2, 5, 5, 5, 5};

np::nurbs_spline2<real> curve(
    controls, weights, knots, 3, false, 1e-10);
const point start = curve.get_start();
const point end = curve.get_end();
const point middle = curve.point_at(2.5);
\end{lstlisting}

\begin{figure}[tbp]
\centering
\begin{tikzpicture}
\begin{axis}[
 width=.91\linewidth,height=6.2cm,axis equal image,
 axis lines=left,xlabel={$x$},ylabel={$y$},
 xmin=-.2,xmax=4.3,ymin=-1.65,ymax=2.6,clip=false,
 grid=major,grid style={manualmuted!15},
 legend style={draw=none,fill=white,font=\small,
   at={(.5,-.22)},anchor=north},legend columns=1]
\addplot[manualmuted,dashed,mark=*,mark size=2pt]
 table[x=x,y=y]{pictures/spline_controls.dat};
\addlegendentry{control polygon}
\addplot[manualblue,very thick,no marks]
 table[x=x,y=y]{pictures/spline_curve.dat};
\addlegendentry{evaluated rational curve}
\node[anchor=north west,font=\small] at (axis cs:0,-.1) {$P_0$, $w_0=1$};
\node[anchor=south,font=\small] at (axis cs:1,2) {$P_1$, $w_1=1$};
\node[anchor=north,font=\small] at (axis cs:2,-1.12) {$P_2$, $w_2=0.7$};
\node[anchor=south,font=\small] at (axis cs:3,2) {$P_3$, $w_3=1.5$};
\node[anchor=north east,font=\small] at (axis cs:4,-.1) {$P_4$, $w_4=1$};
\end{axis}
\end{tikzpicture}
\caption{Control points steer the curve; they are generally not samples on
it. The blue path is evaluated by the library. The clamped endpoint blocks in
this example make the evaluated endpoints coincide with $P_0$ and $P_4$.}
\label{fig:spline-controls}
\end{figure}

The complete-definition constructors accept three \api{std::vector}s followed
by the degree. Omitting the closure argument constructs an open curve; an
optional tolerance follows the degree in that overload. Supplying
\api{closed} before the tolerance enables the seam requirement. In 3D, replace
\api{point2} and \api{nurbs_spline2} with their 3D counterparts. There is no
default empty spline object: construction establishes a usable definition.

An open curve is allowed to have coincident endpoints. A closed curve requires
that coincidence, but does not by itself guarantee a smooth seam, matching
tangent vectors, matching curvature, or a periodic knot representation. For
the ordinary definition check, endpoint separation must not exceed
$16\,\tau\,\max(1,|\hbox{control coordinates}|)$, where $\tau$ is the definition
tolerance. Closed interpolation also checks its first and last input samples
against the supplied tolerance before solving.

The getter families \api{get_control_points()}, \api{get_weights()}, and
\api{get_knots()} return constant references for ordinary splines.
\api{get_control_point(i)}, \api{get_weight(i)}, and \api{get_knot(i)} select
individual values; an invalid index throws \api{std::out_of_range}.
\api{degree()}, \api{tolerance()}, and \api{is_closed()} report metadata.
\api{get_start()} and \api{get_end()} return constant references to evaluated,
cached endpoints, not editable definition fields.

\subsection{Starting with a standard straight curve}

When only endpoints and a number of controls are known, the standard
constructor lays the control points uniformly along the segment using
\api{lerp}, assigns unit weights, and builds standard open-clamped knots. Its
native domain is $[0,\operatorname{distance}(\mathrm{start},\mathrm{end})]$.
The degree must be positive and smaller than the requested control count.
The endpoints must define a finite, positive domain length.

\begin{lstlisting}[language=C++,caption={Standard construction and knot replacement.}]
np::nurbs_spline3<real> straight(
    np::point3<real>{0, 0, 0},
    np::point3<real>{4, 0, 3}, 6, 3);
// Native domain is initially [0, 5].
straight.set_standard_knots(10.0, 20.0);
// Explicit-tolerance overload, without a closure argument:
np::nurbs_spline2<real> straight2(
    point{0, 0}, point{4, 0}, 5, 2, 1e-10);
straight2.set_standard_knots(8.0);  // Domain [0, 8].
\end{lstlisting}

\api{set_standard_knots(s_end)} uses $[0,s_\mathrm{end}]$.
\api{set_standard_knots(s_start,s_end)} uses the stated bounds. Both repeat
each endpoint $p+1$ times and distribute any simple interior knots uniformly.
If $C=p+1$, the two endpoint blocks form the entire knot vector. The generated
spans must remain distinct and have finite positive lengths with finite
reciprocals in \api{REAL}. Invalid domains or numerically unusable spans are
rejected before the curve changes.

Replacing standard knots can change the distribution of points along an
existing curve. If the intent is only to rescale its parameter, apply one
affine transformation to \emph{every} old knot and pass the complete result
to \api{set_knots}; this preserves the geometric locus while changing its
native $s$ values and derivative scale. A closure flag on the straight
constructor still requires the ordinary endpoint check; it does not invent a
loop between two different endpoints.

\begin{figure}[tbp]
\centering
\begin{tikzpicture}[x=1.65cm,y=1cm,font=\small]
\draw[->,manualmuted] (-.3,0)--(5.5,0) node[right] {$s$};
\draw[manualblue,line width=3pt] (0,0)--(5,0);
\foreach \x/\lab in {0/{K_0=K_1=K_2=K_3},2/{K_4},5/{K_5=K_6=K_7=K_8}}{
 \draw[manualink] (\x,-.13)--(\x,.13);
 \node[above=6pt] at (\x,0) {$\lab$};}
\node[below=6pt] at (0,0) {$s_{\min}=0$};
\node[below=6pt] at (2,0) {$2$};
\node[below=6pt] at (5,0) {$s_{\max}=5$};
\draw[<->,manualteal] (0,-1.05)--(2,-1.05)
 node[midway,below] {first nonempty span};
\draw[<->,manualteal] (2,-1.05)--(5,-1.05)
 node[midway,below] {second nonempty span};
\node[align=center,text=manualmuted] at (2.5,-2.1)
 {Cubic running example: $C=5$, $p=3$, nine stored knots.\\
 The two spans have different native parameter lengths.};
\end{tikzpicture}
\caption{Stored knots and the active domain. Endpoint multiplicities describe
clamping; they are not extra finite intervals to sample.}
\label{fig:spline-knots}
\end{figure}

\section{Weights and exact conic geometry}

For basis functions $N_{i,p}$, a rational curve is
\begin{equation}
 \mathbf C(s)=\frac{\sum_{i=0}^{C-1}N_{i,p}(s)w_i\mathbf P_i}
                         {\sum_{i=0}^{C-1}N_{i,p}(s)w_i}.
\label{eq:rational-curve}
\end{equation}
The library evaluates only the nonzero basis functions in the selected knot
span. With positive weights, a control point with a larger relative weight
has more influence where its basis function is nonzero. Uniformly multiplying
the weights leaves the mathematical quotient unchanged, but a very small
common scale can trigger the implementation's near-zero homogeneous-weight
check. In particular, positivity alone is not a guarantee that every
arbitrarily scaled floating-point definition can be evaluated.

A rational quadratic represents a quarter of the unit circle exactly in real
arithmetic. It uses three controls, with the corner control carrying weight
$1/\sqrt2$. The native parameter below is $s\in[0,1]$; it is neither the circle
angle $u$ nor the traveled arc length.

\begin{lstlisting}[language=C++,caption={An exact rational quadratic quarter circle.}]
np::nurbs_spline2<real> quarter(
    std::vector<point>{{1, 0}, {1, 1}, {0, 1}},
    std::vector<real>{1, std::sqrt(0.5), 1},
    std::vector<real>{0, 0, 0, 1, 1, 1}, 2);
const point diagonal = quarter.evaluate(0.5);
const real quarter_curvature = quarter.curvature(0.5);
const auto rational_third = quarter.third_derivative(0.5);
\end{lstlisting}

\begin{figure}[tbp]
\centering
\begin{tikzpicture}
\begin{axis}[
 width=.7\linewidth,height=6.2cm,axis equal image,
 axis lines=left,xlabel={$x$},ylabel={$y$},
 xmin=-.12,xmax=1.25,ymin=-.12,ymax=1.3,
 xtick={0,.5,1},ytick={0,.5,1},grid=major,
 grid style={manualmuted!15}]
\addplot[manualmuted,dashed,mark=*,mark size=2pt]
 coordinates {(1,0) (1,1) (0,1)};
\addplot[manualblue,very thick,no marks]
 table[x=x,y=y]{pictures/rational_circle.dat};
\draw[manualteal,->] (axis cs:0,0)--(axis cs:.7071,.7071);
\addplot[only marks,manualorange,mark=*,mark size=2.5pt]
 coordinates {(.7071,.7071)};
\node[anchor=south west,font=\small] at (axis cs:.72,.73)
 {$s=\frac12$};
\node[anchor=south east,font=\small] at (axis cs:1,1)
 {$w_1=1/\sqrt2$};
\node[anchor=north,font=\small] at (axis cs:.4,.18) {$r=1$};
\end{axis}
\end{tikzpicture}
\caption{A quadratic rational curve lies on the unit circle, even though its
control polygon follows the two sides of a square. Its curvature is one in
inverse world units. The plotted blue points are library evaluations.}
\label{fig:rational-circle}
\end{figure}

Rounded output at the midpoint is
\begin{lstlisting}[caption={Rounded numerical preview; scientific precision is four decimals.}]
position: 7.0711e-01  7.0711e-01
curvature: 1.0000e+00
\end{lstlisting}
Although the homogeneous numerator and denominator have degree two within a
span, their quotient can have a nonzero third derivative. Calling
\api{third_derivative} on this quadratic is therefore meaningful; setting it
to zero solely because $p=2$ would be incorrect.

\section{Evaluation and analytic differential geometry}
\label{sec:spline-derivatives}

\api{evaluate(s)} and \api{point_at(s)} are equivalent position calls.
\api{derivatives_at(s)} computes position, first derivative, and second
derivative together. The result is \api{spline_derivatives2<REAL>} or
\api{spline_derivatives3<REAL>}, with members \api{point}, \api{first}, and
\api{second}. Separate calls are available as \api{first_derivative(s)},
\api{second_derivative(s)}, and \api{third_derivative(s)}.

\begin{lstlisting}[language=C++,caption={Position, derivatives, tangent, and curvature.}]
const real query_s = 2.5;
const auto jet = curve.derivatives_at(query_s);
const auto velocity = jet.first;
const auto acceleration = jet.second;
const auto third = curve.third_derivative(query_s);
const auto unit_tangent = curve.tangent(query_s);
const real curvature = curve.curvature(query_s);
// A consumer chooses how long a derivative arrow should appear:
const point arrow_tip = jet.point + 0.25 * velocity;
\end{lstlisting}

Let $\mathbf A^{(k)}=\sum_i N_{i,p}^{(k)}w_i\mathbf P_i$ and
$W^{(k)}=\sum_i N_{i,p}^{(k)}w_i$. Analytic differentiation is performed in
homogeneous form and then converted with the quotient recurrence
\begin{align}
\mathbf C^{(0)}&=\mathbf A^{(0)}/W^{(0)},\\
\mathbf C^{(r)}&=\frac{1}{W^{(0)}}\left(
 \mathbf A^{(r)}-\sum_{j=1}^{r}\binom rj
 W^{(j)}\mathbf C^{(r-j)}\right),\qquad r=1,2,3.
\label{eq:spline-quotient}
\end{align}
Homogeneous derivatives above the degree are zero, while the quotient
recurrence continues through the requested order. These are analytic basis
derivatives, rather than finite differences between neighboring points.

The derivatives are with respect to the native $s$. If $s$ has units of length,
the first derivative is dimensionless, the second has inverse-length units,
and the third has inverse-square-length units. If $s$ is dimensionless,
all three retain world-length units. More generally,
$\mathbf C^{(r)}$ has units of world length divided by parameter units to
power $r$. Scaling a derivative arrow for a diagram does not change these
API meanings.

\begin{figure}[tbp]
\centering
\begin{tikzpicture}
\begin{axis}[
 width=.47\linewidth,height=5.5cm,axis lines=left,
 xlabel={$s$},ylabel={first derivative},xmin=0,xmax=5,
 grid=major,grid style={manualmuted!15},
 legend style={draw=none,font=\small},legend pos=north east]
\addplot[manualblue,thick] table[x=s,y=dx]{pictures/spline_curve.dat};
\addlegendentry{$C'_x$}
\addplot[manualorange,thick] table[x=s,y=dy]{pictures/spline_curve.dat};
\addlegendentry{$C'_y$}
\end{axis}
\end{tikzpicture}\hfill
\begin{tikzpicture}
\begin{axis}[
 width=.47\linewidth,height=5.5cm,axis lines=left,
 xlabel={$s$},ylabel={$\kappa(s)$},xmin=0,xmax=5,
 grid=major,grid style={manualmuted!15}]
\addplot[manualteal,thick] table[x=s,y=curvature]{pictures/spline_curve.dat};
\end{axis}
\end{tikzpicture}
\caption{Analytic first-derivative components and curvature for the running
cubic curve. Speed and turning vary through the native domain. The plotted
values come from \api{first_derivative} and \api{curvature}.}
\label{fig:spline-differential}
\end{figure}

\subsection{Tangent and curvature contracts}

\api{tangent(s,tolerance)} normalizes the first derivative and throws
\api{std::domain_error} when its length is too small. Its default threshold
is the matching vector type's \api{default_tolerance()}. This is a derivative
length threshold; it is distinct from the spline definition tolerance and
from the world-distance tolerances used by geometric queries.

\api{curvature(s,tangent_tolerance)} returns an unsigned magnitude:
\begin{equation}
 \kappa_2=\frac{|\mathbf C'\mathbin{\times}\mathbf C''|}
                  {\|\mathbf C'\|^3},\qquad
 \kappa_3=\frac{\|\mathbf C'\mathbin{\times}\mathbf C''\|}
                  {\|\mathbf C'\|^3}.
\end{equation}
The 2D cross product in this formula is a scalar. Curvature has inverse
world-length units and is independent of regular reparameterization while
the required arithmetic remains representable. It does not encode the sign
of left or right turning. The threshold must be finite and nonnegative;
an invalid threshold throws \api{std::invalid_argument}. A nonfinite or
insufficient first-derivative length, a near-zero homogeneous weight, or an
unrepresentable curvature throws \api{std::domain_error}.

\subsection{Domain edges and repeated knots}

A finite $s$ just outside the active domain by no more than the definition
tolerance is clamped to the nearest bound. A nonfinite $s$, or a larger
excursion, throws \api{std::out_of_range}. This tolerance is a boundary
allowance, not an instruction to extrapolate the spline.

At an internal knot, a derivative that is discontinuous uses the right-hand
nonempty span. At \api{s_min()}, evaluation chooses the first nonempty span
on the right, and at \api{s_max()} it chooses the last nonempty span on the
left. Repeated boundary knots can leave empty spans; those spans are skipped.
This convention is especially relevant when plotting second or third
derivatives near knots with reduced continuity. A single returned derivative
at such a knot must not be mistaken for a two-sided smooth derivative.

\section{Global interpolation through sample positions}
\label{sec:spline-interpolation}

Use the static \api{interpolate(samples,stations,...)} factory when positions
are measured points that the result should pass through. The positions are
\emph{samples}, not control points. Interpolation solves a global linear
system for a new set of control points and assigns unit weights. It produces
a valid non-rational subset of the NURBS representation, with averaged
clamped knots.

\begin{lstlisting}[language=C++,caption={Interpolation preserves the supplied station domain.}]
const std::vector<point> samples{
    {0, 0}, {1, 1.2}, {2, 0.2}, {3, 1.5}, {4, 1}};
const std::vector<real> stations{2, 3.2, 4.7, 6.5, 8};
auto fit = np::nurbs_spline2<real>::interpolate(
    samples, stations, 3, false, 1e-10);

for (std::size_t i = 0; i < samples.size(); ++i) {
    const point fitted_sample = fit.evaluate(stations[i]);
    const real residual = np::distance(fitted_sample, samples[i]);
    // Inspect residual in the same 2D world as the input.
    static_cast<void>(residual);
}
// fit.s_min() == 2; fit.s_max() == 8.
\end{lstlisting}

\begin{figure}[tbp]
\centering
\begin{tikzpicture}
\begin{axis}[
 width=.9\linewidth,height=6.2cm,axis lines=left,
 xlabel={$x$},ylabel={$y$},xmin=-.25,xmax=4.3,ymin=-.25,ymax=1.85,
 clip=false,
 grid=major,grid style={manualmuted!15},
 legend style={draw=none,fill=white,font=\small},legend pos=south east]
\addplot[manualblue,very thick,no marks]
 table[x=x,y=y]{pictures/interpolation_curve.dat};
\addlegendentry{global cubic interpolant}
\addplot[only marks,manualorange,mark=*,mark size=2.8pt]
 table[x=x,y=y]{pictures/interpolation_samples.dat};
\addlegendentry{input samples}
\node[anchor=north west,font=\small] at (axis cs:0,0) {$s=2$};
\node[anchor=south,font=\small] at (axis cs:1,1.2) {$s=3.2$};
\node[anchor=north,font=\small] at (axis cs:2,.2) {$s=4.7$};
\node[anchor=south,font=\small] at (axis cs:3,1.5) {$s=6.5$};
\node[anchor=south east,font=\small,xshift=-2pt,yshift=2pt]
 at (axis cs:4,1) {$s=8$};
\end{axis}
\end{tikzpicture}
\caption{A curve through measured positions. Each annotation is the original
native station, not a newly normalized coordinate. The library supplies the
evaluated curve and sample positions used in this TikZ picture.}
\label{fig:spline-interpolation}
\end{figure}

The input vectors must have equal size and contain at least two samples.
All samples and stations must be finite; station values must be strictly
increasing. The requested degree must be positive, and the actual degree is
\begin{equation}
 p=\min(\text{requested degree},\text{sample count}-1).
\end{equation}
The default requested degree is three and the default interpolation tolerance
is $10^{-10}$. Invalid samples, stations, closure, or tolerance throw
\api{std::invalid_argument}; a singular interpolation system throws
\api{std::domain_error}. Floating-point residuals should be assessed with a
suitable tolerance rather than exact coordinate equality.

Let the sample stations be $t_0,\ldots,t_n$. The clamped endpoint blocks contain
$p+1$ copies of $t_0$ and $t_n$. The simple interior knots are averages,
\begin{equation}
 K_{j+p}=\frac{t_j+t_{j+1}+\cdots+t_{j+p-1}}{p},
 \qquad j=1,\ldots,n-p.
\end{equation}
For the example, the sole interior knot is
$(3.2+4.7+6.5)/3=4.8$. Solving
$\sum_i N_{i,p}(t_j)\mathbf P_i=\mathbf Q_j$ independently for each coordinate
makes all samples influence the global result. Adding or moving one sample
can therefore change more than the adjacent visible arc.

\begin{manualnote}[Stations remain stations]
The caller supplies the arc-length stations; interpolation does not compute
them from the sample positions. Their values are retained, without silently
normalizing the domain. Even if the stations describe measured arc length,
the interpolated curve between stations does not promise unit speed or an
exact continuous arc-length parameterization.
\end{manualnote}

\subsection{Closed input and replacing an existing curve}

The overload with \api{closed=true} requires coincident first and final sample
positions at \emph{different}, strictly increasing station values. Repeat
the closing position, not its station. This constructs a closed interpolant
using the same averaged clamped-knot approach. It does not impose periodic
derivatives at the seam.

\begin{lstlisting}[language=C++,caption={Closed interpolation and adopting a new sample set.}]
const std::vector<point> loop_samples{
    {0, 0}, {1, 0}, {1, 1}, {0, 1}, {0, 0}};
const std::vector<real> loop_stations{0, 1, 2, 3, 4};
auto loop = np::nurbs_spline2<real>::interpolate(
    loop_samples, loop_stations, 3, true, 1e-10);

auto refitted = curve;
refitted.adopt_to_points(samples, stations, 2, false, 1e-10);
// Unlike definition setters, adoption creates a new degree-two fit.
\end{lstlisting}

\api{adopt_to_points} has the same open and closed overload families as
\api{interpolate}. It first constructs the interpolant and then replaces the
current object. A failed solve leaves the original object available. Its
requested degree, tolerance, and closure describe the replacement; they are
not constrained to the old curve's metadata. Both array-backed curve types
also supply these interpolation and adoption operations in their own world.

\section{Two different notions of length}

\api{get_polygon_length()} sums distances between consecutive stored control
points:
\begin{equation}
 L_\mathrm{polygon}=\sum_{i=1}^{C-1}\|\mathbf P_i-\mathbf P_{i-1}\|.
\end{equation}
This is the length of the \emph{control polygon}, not an integration of the
curve. It depends on neither weights nor knots. Even for a closed curve, it
does not add an extra last-to-first edge. For ordinary splines this getter
requires no evaluation; for array-backed splines it first validates the
current public definition.

\api{approximate_arc_length(segment_count)} instead evaluates a polyline at
uniformly spaced native parameters, including both active-domain endpoints:
\begin{equation}
 s_j=s_\mathrm{min}+\frac{j}{M}(s_\mathrm{max}-s_\mathrm{min}),\qquad
 L_M=\sum_{j=1}^{M}\|\mathbf C(s_j)-\mathbf C(s_{j-1})\|.
\end{equation}
The default is $M=512$; zero segments throw \api{std::invalid_argument}.
The result is a sampling estimate in world units. It is not an adaptive
quadrature result and provides no requested-error certificate. Tight turns
or short features may need many more segments, especially when the native
parameter allocates little interval length to them.

\begin{lstlisting}[language=C++,caption={Compare a control polygon with sampled curve lengths.}]
const real polygon_length = curve.get_polygon_length();
const real coarse_length = curve.approximate_arc_length(16);
const real finer_length = curve.approximate_arc_length(512);
const real circle_length = quarter.approximate_arc_length(512);
// circle_length approaches pi/2 as the segments resolve the arc.
\end{lstlisting}

\begin{figure}[tbp]
\centering
\begin{tikzpicture}
\begin{axis}[
 width=.87\linewidth,height=5.8cm,axis lines=left,
 xlabel={$x$},ylabel={$y$},xmin=-.15,xmax=4.15,ymin=-.15,ymax=1.45,
 grid=major,grid style={manualmuted!15},
 legend style={draw=none,fill=white,font=\small,
   at={(.5,-.22)},anchor=north},legend columns=1]
\addplot[manualblue,line width=1.5pt]
 table[x=x,y=y]{pictures/spline_curve.dat};
\addlegendentry{densely evaluated curve}
\addplot[manualorange,thick,mark=*,mark size=2pt,each nth point=40]
 table[x=x,y=y]{pictures/spline_curve.dat};
\addlegendentry{coarse uniform-$s$ chords}
\end{axis}
\end{tikzpicture}
\caption{Arc-length sampling joins evaluated points with straight chords.
Uniform native-parameter steps generally produce unequal world-space chord
lengths. The coarse orange polyline illustrates the approximation, not an
adaptive accuracy guarantee.}
\label{fig:spline-length}
\end{figure}

When comparing successive estimates, retain the sampling count along with
the value. Agreement at two finite counts is useful evidence for that curve,
but cannot establish that all continuous features were resolved. The same
qualification applies to the sample-based geometric queries in
Chapter~\ref{chap:algorithms}.

\section{Checked editing and flattened scalar parameters}
\label{sec:spline-editing}

Ordinary spline definitions are edited through setters. The single-value
setters validate a complete detached candidate: \api{set_control_point}
edits one position, \api{set_weight} one weight, and \api{set_knot} one knot.
The vector replacements \api{set_control_points},
\api{set_weights}, and \api{set_knots} retain the other definition fields.
\api{set_tolerance} and \api{set_closed} also recheck the complete candidate,
including any seam requirement. None of these operations changes the degree.

\api{set_definition(control_points,weights,knots,closed,tolerance)} replaces
every mutable definition field together while retaining the current degree.
It is the appropriate call for changing correlated counts, rescaling an
entire knot vector, or editing both sides of a closed seam. Successful edits
refresh the cached endpoints; failed edits preserve the original curve.
References, iterators, and pointers previously obtained from ordinary
definition getters become invalid after a successful setter.

\begin{lstlisting}[language=C++,caption={Atomic bulk edits keep correlated changes together.}]
auto edited = curve;
auto edited_controls = edited.get_control_points();  // Detached copy.
auto edited_knots = edited.get_knots();              // Detached copy.
edited_controls.front() = {-1, 0};
edited_controls.back() = {5, 0};
for (real& knot : edited_knots) {
    knot = 10 + 2 * knot;
}
edited.set_definition(
    edited_controls, edited.get_weights(), edited_knots,
    edited.is_closed(), edited.tolerance());
// New endpoints and domain [10, 20] are immediately available.
\end{lstlisting}

\subsection{An indexable view for fitting code}

\api{number_of_parameters()}, \api{get_parameter(index)}, and
\api{parameter_name(index)} expose the same flattened scalar ordering in all
four spline types. Control coordinates come first in point-major order,
followed by all weights, followed by all knots. In dimension $d$, the total is
$(d+2)C+p+1$ scalars. Degree, tolerance, and closure are deliberately outside
this scalar sequence. Names are \api{P[0].x}, \api{P[0].y}, optionally
\api{P[0].z}, then later controls, \api{W[0]}, and finally \api{K[0]}, all
with zero-based indices.

\begin{figure}[tbp]
\centering
\begin{tikzpicture}[x=.54cm,y=1cm,font=\small]
\fill[manualblue!13] (0,0) rectangle (6,1);
\fill[manualteal!15] (6,0) rectangle (9,1);
\fill[manualorange!18] (9,0) rectangle (15,1);
\draw[manualmuted] (0,0) rectangle (15,1);
\foreach \x in {1,...,14} \draw[manualmuted!45] (\x,0)--(\x,1);
\foreach \x/\lab in {0/{x_0},1/{y_0},2/{x_1},3/{y_1},4/{x_2},5/{y_2},6/{w_0},7/{w_1},8/{w_2},9/{K_0},10/{K_1},11/{K_2},12/{K_3},13/{K_4},14/{K_5}}
 \node at (\x+.5,.5) {$\lab$};
\draw[<->,manualblue] (0,-.4)--(6,-.4)
 node[midway,below,align=center] {six\\coordinates};
\draw[<->,manualteal] (6,-.4)--(9,-.4)
 node[midway,below,align=center] {three\\weights};
\draw[<->,manualorange] (9,-.4)--(15,-.4)
 node[midway,below,align=center] {six\\knots};
\node[above,align=center] at (7.5,1.25)
 {2D quadratic example: $C=3$, $p=2$, $N=15$.};
\node[below,align=center,text=manualmuted] at (7.5,-1.4)
 {Scalar indices: coordinates $[0,6)$, weights $[6,9)$, knots $[9,15)$.};
\end{tikzpicture}
\caption{Flattened numeric ordering for an ordinary spline's indexed
parameters and an array-backed spline's public \api{parameters} storage.
Typed metadata is stored separately.}
\label{fig:spline-layout}
\end{figure}

\api{set_parameter(index,value)} returns \api{true} only if the full edited
candidate is accepted. It returns \api{false} for an invalid index or a rejected
definition, retaining the original state. Allocation failures can still
throw. This return-value convention differs from the named component setters,
whose invalid candidate or index is reported by an exception.

\api{set_parameters(const REAL*)} copies exactly
\api{number_of_parameters()} values in this ordering before validating.
The pointer is not retained. It must be nonnull and address a sufficiently
long readable buffer; the library cannot discover a raw pointer's extent.
For a complete numeric change, prepare a detached vector and call once:

\begin{lstlisting}[language=C++,caption={Indexed parameter names and one validated bulk copy.}]
std::vector<real> candidate(curve.number_of_parameters());
for (std::size_t i = 0; i < candidate.size(); ++i) {
    candidate[i] = curve.get_parameter(i);
    const std::string name = curve.parameter_name(i);
    static_cast<void>(name);
}
candidate[3] += 0.1;  // P[1].y in the 2D point-major layout.
auto scalar_edited = curve;
scalar_edited.set_parameters(candidate.data());
const bool accepted = scalar_edited.set_parameter(0, -0.2);
static_cast<void>(accepted);
\end{lstlisting}

There is no endpoint setter: an endpoint is derived from the definition.
Change the appropriate control-point definition, then read the refreshed
endpoint. Do not assume the first control point is always the active start:
that familiar identity depends on clamped endpoint knots.

\section{Array-backed curves for direct numeric work}
\label{sec:array-splines}

The array-backed classes are independent types, not subclasses of the
ordinary splines. They offer the same core geometry, checked setters, and
interpolation calls, plus a public \api{parameters} array and layout/range
queries. Their aliases \api{real_t} and \api{parameter_index_range} name the
scalar type and the allocation-free range of flattened scalar indices.

Construction can use endpoints with a control count and degree, component
vectors with a degree, a flattened \api{std::valarray<REAL>} with a degree,
or an explicit copy from the corresponding ordinary spline. The optional
closure and tolerance forms mirror these choices. The flattened constructor
derives $C$ from the scalar count $N$:
\begin{equation}
 \begin{array}{ll}
 \text{2D:}&N=4C+p+1,\\
 \text{3D:}&N=5C+p+1.
 \end{array}
\end{equation}
The inferred count must be integral and exceed the degree. A 2D sequence
cannot be interpreted as a 3D sequence simply by choosing another type.

\begin{lstlisting}[language=C++,caption={A flattened quadratic and same-world conversions.}]
const std::valarray<real> flat_quarter{
    1, 0,  1, 1,  0, 1,       // Three XY control points.
    1, std::sqrt(0.5), 1,       // Three weights.
    0, 0, 0, 1, 1, 1};        // Six knots.
np::nurbs_arr_spline2<real> array_quarter(flat_quarter, 2);

np::nurbs_arr_spline2<real> array_curve(curve);
np::nurbs_spline2<real> vector_curve(array_curve);
auto vector_copy = array_curve.to_nurbs_spline();
const point array_point = array_curve.evaluate(2.5);
\end{lstlisting}

These conversions copy the complete definition and preserve control points,
weights, knots, degree, tolerance, closure, and native domain. They validate
the current source array and can fail if it was edited into an invalid state.
The same operations exist between \api{nurbs_arr_spline3} and
\api{nurbs_spline3}. For shared ownership, the matching
\api{make_nurbs_spline2}, \api{make_nurbs_spline3},
\api{make_nurbs_arr_spline2}, and \api{make_nurbs_arr_spline3} allocation helpers
are covered in Chapter~\ref{chap:creators}. Ordinary allocation helpers also
accept component valarrays; that convenience overload does not turn an
ordinary curve into the flattened array-backed type.

\subsection{Direct editing, validation, and range meanings}

Direct array edits intentionally bypass checked setters. They can temporarily
leave invalid weights, bad knot order, a broken closed seam, or a malformed
array size. \api{validate()} checks the current layout, values, native domain,
and required seam. Geometric operations perform their own current-state
validation before using the array; they reject malformed storage before
indexing through its inferred layout.

No endpoint or other geometry cache is retained. \api{get_start()},
\api{get_end()}, \api{s_min()}, and \api{s_max()} derive their answers from
the current validated storage. Point getters reconstruct point values from
their coordinates; they do not alias adjacent scalars as a point object.
\api{get_control_points()}, \api{get_weights()}, and \api{get_knots()} return
detached vectors, unlike the ordinary spline getters' constant references.
Layout and component-access getters check shape; call \api{validate()} when
the whole definition, including numerical values and closure, must be checked.
\api{number_of_parameters()} reports the public array size even if its current
shape is malformed.

\api{control_point_count()}, \api{weight_count()}, and \api{knot_count()}
report component counts. The three range getters return \emph{flattened
scalar indices}, despite the word \api{count} in their names:

\begin{center}
\small
\begin{tabularx}{\linewidth}{@{}lXX@{}}
\toprule
Range getter & 2D indices & 3D indices\\
\midrule
\api{get_control_point_count_range()} & $[0,2C)$ & $[0,3C)$\\
\api{get_weight_count_range()} & $[2C,3C)$ & $[3C,4C)$\\
\api{get_knot_count_range()} & $[3C,N)$ & $[4C,N)$\\
\bottomrule
\end{tabularx}
\end{center}

\begin{lstlisting}[language=C++,caption={Use scalar ranges without assuming their offsets.}]
auto optimized = array_curve;
for (const auto i : optimized.get_weight_count_range()) {
    optimized.parameters[i] *= 1.25;
}
optimized.validate();

auto proposed = optimized.parameters;  // Detached valarray copy.
proposed[3] += 0.2;                    // Modify P[1].y.
optimized.set_parameters(proposed);    // Check before committing.
\end{lstlisting}

A range captures bounds when requested. Obtain a new range after replacing
or resizing \api{parameters}; do not reuse stale bounds. To preserve atomic
updates, prefer the array-specific
\api{set_parameters(std::valarray<REAL>)} or the common checked setter
families. The raw-pointer overload copies the current array's scalar count
and supports a source pointer into that current storage by detaching first.
\api{set_parameter} additionally returns \api{false} for a malformed current
shape, while allocation failure can still throw.

\subsection{Choosing the storage contract}

Use an ordinary spline when callers should see a continuously valid object
with private storage, stable metadata access, and cached endpoint queries.
Use an array-backed spline when numeric fitting code needs an openly
replaceable scalar array and can take responsibility for validating its
edits. Both evaluate in their own Cartesian world and use analytic
derivatives. Neither conversion nor direct array arithmetic is a coordinate
transform.

The persistence, projection, drawing, distance, and intersection interfaces
are documented with their accepted types in later chapters and the API
reference. When an operation accepts an ordinary spline specifically, call
\api{to_nurbs_spline()} explicitly; do not assume inheritance or an implicit
conversion. The ordinary 3D spline supplies \api{tag_write} and its static
\api{tag_read}; array-backed splines can use an explicit ordinary conversion
for that 3D tagged representation.
